\documentclass[11pt,a4paper]{article}
\usepackage[T1]{fontenc}
\usepackage{lmodern,amsmath,amssymb,amsthm,mathtools,booktabs,microtype}
\usepackage[margin=26mm]{geometry}
\usepackage[colorlinks=true,linkcolor=blue,citecolor=blue,urlcolor=blue]{hyperref}
\setlength{\emergencystretch}{3em}
\newtheorem{theorem}{Theorem}[section]\newtheorem{lemma}[theorem]{Lemma}
\newtheorem{proposition}[theorem]{Proposition}\newtheorem{corollary}[theorem]{Corollary}
\theoremstyle{definition}\newtheorem{definition}[theorem]{Definition}
\theoremstyle{remark}\newtheorem{remark}[theorem]{Remark}
\newcommand{\Q}{\mathbb Q}\newcommand{\Z}{\mathbb Z}
\DeclareMathOperator{\sfree}{sf}
\title{Common twists for several pole inputs}
\author{}\date{14 November 2026}
\begin{document}\maketitle
\begin{abstract}
We apply the finite-family prescribed-class twist theorem in the supplied
row-5a audit to the congruent-number pole equations. Inputs with a common
nontrivial input square class and $k$ distinct nontrivial complement square
classes can use $k+1$ retained rational parameters under an explicit common-twist
hypothesis. Pairwise-coprime admissible complements satisfy that hypothesis
by the supplied theorem. We also prove a change of base square class which
extends the application to some overlapping families; two explicit prepared
inputs have complements sharing the prime 89. The equations and witness maps
are exact, including trivial complement classes and repeated classes.
These are reductions of pole subsystems. They do not prove a smaller complete
25- or 13-variable definition, whose present witnesses have different input
square classes.
\end{abstract}

\section{Source theorem and scope}
Write $C_n:y^2=x^3-n^2x$ for $n>0$ rational. For a positive rational
number, its squarefree representative is the unique positive squarefree
integer in its rational square class. Two positive squarefree integers
$\beta,\theta$ are called admissible here when
\[
 \beta>1,\quad\beta\equiv\theta\equiv1\pmod8,\quad
 \gcd(\beta,\theta)=1,
\]
and $\theta$ is a quadratic residue at every prime dividing $\beta$,
while $\beta$ is a quadratic residue at every prime dividing $\theta$.
The conditions on the latter set are vacuous for $\theta=1$.

We use precisely the following consequence of the theorem entitled
\emph{Simultaneous prescribed-class twists} in the supplied audit
\cite{Audit}.
To avoid confusing that strengthening with the original two-twist theorem,
we state it explicitly as an input.

\begin{theorem}[Retained finite-family input]\label{thm:input}
Let $\beta>1$ be positive squarefree and $1\pmod8$. Let
$\eta_1,\ldots,\eta_s$ be pairwise coprime positive squarefree
integers, each admissible with $\beta$. There is a positive squarefree
$m$ such that $C_m$ and all $C_{m\eta_j}$ have infinite-order rational
points with abscissa in $\beta\Q^{\times2}$.
\end{theorem}

This is a weaker conclusion than the supplied audit's theorem, which also
states rank, Selmer and density conclusions. Those stronger statements
are not needed here. The finite-family theorem is the supplied audit's
strengthening of the congruent-number construction in Jaswal's Hexagon
preprint \cite{Jaswal}, not a conclusion of separate applications of
Jaswal's two-twist Theorem~1.2. We retain its stated hypotheses and do not
claim to have newly audited its external analytic dependencies. The supplied
TeX source is preserved unchanged under
\nolinkurl{sources/hexagon}; its hash is bound in the exact certificate.

A positive rational $b$ is prepared for the pole equations of
\cite{Congruent} when
\[
 0<b<1,\quad b\text{ has square reduced denominator},\quad
 b\in\Q_2^{\times2},\quad v_2(b)\ge3.
\]
For such a nonsquare $b$, the decompositions
\[
                       b=\beta z^2,\qquad 1-b=\theta c^2
\]
with positive squarefree representatives have $(\beta,\theta)$
admissible. Briefly, the coprime numerators of $b$ and $1-b$ over their
common square denominator give coprimality and the mutual residue
conditions. The square unit of $b$ at two gives $\beta\equiv1\pmod8$;
its even valuation is at least four, so $1-b\equiv1\pmod{16}$ and
$\theta\equiv1\pmod8$. These statements are proved with the same
conventions in \cite{Congruent}.

\section{The common-twist implication and its exact parameter count}
The pole equations are
\begin{equation}\label{eq:pole}
 e^2+bp^2=1,\qquad uv=e(1-b),\qquad u^2-v^2=bq^2.
\end{equation}
A retained parameter is a freely chosen rational coordinate before the
single finite-algebra elimination. Square roots of specified rational
functions are finite choices; their rationality must still be imposed.
We use the one-coordinate encoding in \cite{Joint}, which can be shared
with the existing finite choices in a larger construction.

\begin{theorem}\label{thm:common}
Let $b_1,\ldots,b_r$ be prepared nonsquare inputs with the same
squarefree representative $\beta>1$. Write
\[
 b_i=\beta z_i^2,\qquad 1-b_i=\theta_i c_i^2,
 \qquad z_i,c_i\in\Q^\times.
\]
Let $\mathcal T$ be the set of their distinct nontrivial complement
representatives, and put $k=|\mathcal T|$.
Suppose one positive rational $n$ exists for which $C_n$ and every
$C_{n\theta}$, $\theta\in\mathcal T$, have infinite-order rational
points with abscissa in $\beta\Q^{\times2}$.
Then the joint pole systems have witnesses with only the $k+1$ retained
parameters
\[
                              e,\quad (w_\theta)_{\theta\in\mathcal T},
\]
together with finite square roots. They may be chosen with
$0<e<1$ and all $u_i,p_i,q_i$ nonzero. The attainable retained tuples
are Zariski dense in $\mathbb A^{k+1}$.
\end{theorem}
\begin{proof}
Choose a non-torsion rational point $(\beta a_0^2,y_0)$ on $C_n$,
and independently points $(\beta a_\theta^2,y_\theta)$ on
$C_{n\theta}$. Their abscissas are positive and their ordinates
nonzero. Define
\[
 e=\frac{n}{\beta a_0^2},\quad
 p_* =\frac{y_0}{\beta^2a_0^3},\quad
 w_\theta=\frac{a_\theta}{a_0},\quad
 q_\theta=\frac{y_\theta}{\beta^2a_0^3w_\theta^2}.
\]
The first curve equation gives
\begin{equation}\label{eq:finite}
 \beta p_*^2=1-e^2,\qquad
 \beta w_\theta^2q_\theta^2=w_\theta^4-e^2\theta^2.
\end{equation}
The latter equality follows by dividing the corresponding second curve
equation by $\beta^3a_0^6w_\theta^2$. A positive abscissa on $C_n$ is larger than $n$
unless it is two-torsion. Hence $0<e<1$.

For the trivial complement put $w_1=1$ and $q_1=p_*$; these satisfy
\eqref{eq:finite} with $\theta=1$. For each input set
\begin{equation}\label{eq:maps}
 p_i=p_*/z_i,\qquad u_i=c_iw_{\theta_i},\qquad
 v_i=e\theta_i c_i/w_{\theta_i},\qquad
 q_i=c_iq_{\theta_i}/z_i.
\end{equation}
Substitution in \eqref{eq:pole} proves the first equation without a
factor, and the other two after multiplication by $c_i^2$. All divisions
are by nonzero quantities, and all specified nonzero witnesses follow
from the non-torsion ordinates. Repeated complement classes use the same
$w_\theta,q_\theta$, not extra retained coordinates.

Vary positive odd multiples of every chosen point independently. On a
curve $C_n$, the abscissa square-class homomorphism is constant on odd
multiples. This follows from the chord-and-tangent product identity
and the duplication formula, as detailed in \cite{Joint}. These odd
multiples have infinitely many distinct abscissas: equality of two
abscissas means equality of the points up to sign, which would give a
torsion multiple of a non-torsion point. Thus the first point gives
infinitely many $e$. With it fixed, the independently varying second
points give independently infinitely many $w_\theta$. A polynomial
vanishing on all these tuples first vanishes in all the $w_\theta$,
and then each coefficient vanishes at infinitely many $e$. This proves
density, also for $k=0$.

For the parameter count in a \emph{uniform} polynomial construction,
one must not introduce the variable squarefree representatives as
coefficients. Choose an existing input $b_*=b_1$ instead of $\beta$,
and for each nontrivial complement class choose an existing representative
$d_\theta=1-b_{j(\theta)}=\theta c_{j(\theta)}^2$.
Rescale $w_\theta$ to $\widetilde w_\theta=c_{j(\theta)}w_\theta$,
$p_*$ to $\widetilde p=p_*/z_1$, and $q_\theta$ to
$\widetilde q_\theta=c_{j(\theta)}q_\theta/z_1$. Their equations are
\[
 b_*\widetilde p^2=1-e^2,\qquad
 b_*\widetilde w_\theta^2\widetilde q_\theta^2
       =\widetilde w_\theta^4-e^2d_\theta^2.
\]
For input $i$ the remaining finite square roots are
\[
 Z_i^2=b_i/b_*,\qquad C_i^2=(1-b_i)/d_{\theta_i}.
\]
The formulas \eqref{eq:maps} hold with
$(p_*,w_\theta,q_\theta,z_i,c_i,\theta)$ replaced by
$(\widetilde p,\widetilde w_\theta,\widetilde q_\theta,Z_i,C_i,d_\theta)$.
For trivial complement use $d_1=1$, $\widetilde w_1=1$ and
$\widetilde q_1=\widetilde p$. Thus only already present input
functions, not their arithmetically extracted squarefree parts, enter
the coefficient construction. A fixed partition into complement classes
selects the representatives; there are only finitely many such partitions
for a fixed number of inputs, so alternatives may be handled by a finite
disjunction when needed.

All these radicands and their denominators are guarded nonzero. The
separation construction in \cite{Joint} encodes their rational square
roots using one shared additional coordinate. The retained count is
$k+1$; a stand-alone polynomial also needs that separating coordinate,
which a larger construction may already possess. No variable for
$\beta$, $\theta$, or $n$ has been omitted from this accounting: they
are used only in the existence proof, not in the displayed uniform
finite encoding.
\end{proof}

The result is existential equivalence to the collection of individual
pole conditions \emph{on the stated prepared locus and under the common-twist
input}: its witnesses give every individual condition, and the theorem
gives a common witness on that whole locus. It is not a bijection with
all previously chosen witnesses. Further guards may be imposed only if
their polynomials remain nonzero after the displayed substitutions;
density does not imply arbitrary equalities or numerical coincidences.

\begin{corollary}\label{cor:coprime}
If the distinct nontrivial $\theta_i$ are pairwise coprime,
Theorem~\ref{thm:input} supplies the common-twist hypothesis in
Theorem~\ref{thm:common}. In particular two prepared inputs with the
same nontrivial input class and two coprime nontrivial complement
classes need three retained parameters rather than four.
\end{corollary}
\begin{proof}
Apply the finite-family input to the distinct complement classes.
Their admissibility follows from preparation. Then apply
Theorem~\ref{thm:common}.
\end{proof}

\section{A change of base class for overlapping complements}
The pairwise-coprime criterion is sufficient but not necessary. The next
algebraic observation enlarges its scope without strengthening the
finite-family input.

\begin{proposition}\label{prop:change}
Let $\theta_0=1,\theta_1,\ldots,\theta_k$ be positive squarefree
integers. Suppose a positive squarefree $\delta$ exists such that the
distinct nontrivial integers
\[
                  \eta_i=\sfree(\delta\theta_i)\quad(0\le i\le k)
\]
are pairwise coprime and each is admissible with $\beta$.
Then Theorem~\ref{thm:input} supplies a positive rational $n$ such that
$C_n$ and all $C_{n\theta_i}$ have the required infinite-order points
in the abscissa class $\beta$.
\end{proposition}
\begin{proof}
Choose $m$ from that theorem for the distinct nontrivial $\eta_i$,
using its base curve $C_m$ for any $\eta_i=1$. Set $n=m\delta$.
As $\delta,\theta_i$ are squarefree, with
$h_i=\gcd(\delta,\theta_i)$ one has
\[
                  \delta\theta_i=\eta_i h_i^2.
\]
The isomorphism
\[
 C_{m\eta_i}\longrightarrow C_{n\theta_i},\qquad
                   (x,y)\longmapsto(h_i^2x,h_i^3y)
\]
preserves infinite order and multiplies the abscissa by a rational square.
For $i=0$ it supplies $C_n$. The rational parameter $n$ need not itself
be squarefree; neither the pole formulas nor Theorem~\ref{thm:common}
requires this. This is a change of base in the family of curves, not
a claim that arbitrary overlaps satisfy the displayed admissibility
conditions.
\end{proof}

\begin{corollary}\label{cor:star}
Suppose the nontrivial complement classes have the form
$\theta_i=a d_i$, where $a,d_1,\ldots,d_k$ are pairwise coprime
positive squarefree integers, all admissible with $\beta$.
Then Theorem~\ref{thm:common} applies even when $a>1$.
\end{corollary}
\begin{proof}
Use $\delta=a$. The transformed classes are $a,d_1,\ldots,d_k$.
Repeated or trivial classes are removed before applying the input.
\end{proof}

For two complements put $a=\gcd(\theta_1,\theta_2)$ and
$d_i=\theta_i/a$. These three integers are automatically pairwise
coprime. Admissibility of all three must still be checked: admissibility
of their two products alone does not imply it. In particular the residue
of $a$ modulo eight and its individual residue symbols at primes dividing
$\beta$ cannot be dropped.

\section{Explicit prepared inputs and exact tests}
The first pair is
\[
 b_1=\frac{272}{361}=17\left(\frac4{19}\right)^2,\qquad
 b_2=\frac{1088}{1225}=17\left(\frac8{35}\right)^2.
\]
Their complement classes are $89$ and $137$; these are coprime and
admissible with $17$. Both inputs are prepared and fall under
Corollary~\ref{cor:coprime}.

For an overlapping example take
\[
 b_1=\frac{2448}{14641}=17\left(\frac{12}{121}\right)^2,\qquad
 b_2=\frac{4352}{27225}=17\left(\frac{16}{165}\right)^2.
\]
Their complements are
\[
 1-b_1=\frac{89\cdot137}{121^2},\qquad
 1-b_2=\frac{89\cdot257}{165^2}.
\]
The first criteria do not apply directly, since the two representatives
share 89. Instead use $\delta=89$: the transformed classes are
$89,137,257$. Each prime is $1\pmod8$ and a quadratic residue modulo
17; reciprocity also gives the reverse symbols, or they may be checked
directly. Thus Proposition~\ref{prop:change} supplies the common twist.
The two inputs again require three retained parameters.

The exact checker verifies preparation, prime factorization and all
residue symbols, the witness substitutions modulo \eqref{eq:finite},
and the change-of-base curve isomorphism. The examples establish the
hypotheses of the realization theorem; no explicit huge twist or full
numeric pole witness is asserted to have been computed for these pairs.

\section{Relation to the variable-reduction objective}
These results remove the extra unproved multi-twist hypothesis from the
v29 common-first-parameter theorem on the two explicitly described
compatibility loci. They do not force the input classes to coincide.
In particular the current 13-variable completeness witnesses have
\emph{different} first square classes at fresh selected primes, as
proved in \cite{Fresh}. Nor does this argument make the 25-variable
pole inputs compatible. The global counts consequently remain 25 and
13. No new quartic presentation or quartic count follows from the
finite-square-root encoding used here.

The unresolved application is a completeness construction whose pole
inputs have a common first class while retaining the required local
conditions and height inequalities. The present proof supplies an exact
realization theorem once that condition is obtained; it does not assume
that the old infinite-slope construction can impose it.

\begin{thebibliography}{9}
\bibitem{Audit} \emph{A row-5a audit and a simultaneous-twist strengthening:
On Manvir Jaswal's congruent-number route to Diophantine definability},
supplied project audit, 13 November 2026, theorem entitled
\emph{Simultaneous prescribed-class twists} (source label
\texttt{thm:sim}). Unmodified source:
\nolinkurl{sources/hexagon/jaswal_row5a_audit_and_simultaneous_twists_20261113.tex}.
\bibitem{Jaswal} M.~Jaswal, \emph{A proof that $\Z$ is Diophantine in $\Q$
through congruent-number twists}, Hexagon \texttt{2611.00331v2},
12 November 2026, Theorem~1.2.
\url{https://hexagonmath.org/pdf/2611.00331v2}.
\bibitem{Joint} \emph{Joint realization of compatible pole conditions},
project manuscript, 14 November 2026.
\nolinkurl{papers/pole_parameters/joint_pole_realization_20261114_v29.tex}.
\bibitem{Fresh} \emph{Fresh-Prime and Fixed-Scale Obstructions to Joint Pole Realization},
project manuscript, 16 November 2026.
\par
{\raggedright
\nolinkurl{negative_results/papers/pole_parameters/fresh_prime_squareclass_obstruction_20261116_v42.tex}.\par}
\bibitem{Congruent} \emph{Pole equations from simultaneous congruent-number twists},
project manuscript, 14 November 2026.
\nolinkurl{papers/pole_parameters/congruent_twist_poles_20261114_v27.tex}.
\end{thebibliography}
\end{document}
